\documentclass[11pt]{article}

\usepackage[sfdefault]{FiraSans}
\usepackage[T1]{fontenc}
\renewcommand*\oldstylenums[1]{{\firaoldstyle #1}}

\usepackage[english]{babel}
\usepackage{fancyhdr}
\usepackage{booktabs}
\usepackage{hyperref}\hypersetup{colorlinks=true}
\usepackage[dvipsnames]{xcolor}

\usepackage{amsmath,amssymb,amsfonts,textcomp}
\setlength{\tabcolsep}{8pt}
\usepackage{setspace}
\usepackage{graphicx}
\usepackage[margin=1in]{geometry}
\setlength{\parindent}{0pt}
\pagestyle{fancy}

\usepackage{titlesec}
\titleformat{\section}
  {\normalfont\Large\scshape\bfseries}{\thesection}{1em}{}
\titlespacing{\section}{0pt}{10pt}{0pt}
\titleformat{\subsection}
  {\normalfont\bfseries}{\thesection}{1em}{}
\titlespacing{\subsection}{0pt}{6pt}{0pt}

\lhead{ECON3500 -- Econometrics and Applications}
\rhead{In-Class Activity: Instrumenting Education \textit{(Solutions)}}

\newcommand{\ans}[1]{%
  \vspace{2pt}%
  {\color{BrickRed}\textsc{\textbf{Answer.}}} {\itshape #1}\par\vspace{4pt}%
}

\setlength\parskip{0.10in}

\begin{document}
\thispagestyle{plain}
\singlespacing

ECON3500 Econometrics and Applications \hfill Spring 2026

\begin{center}
\Large{\textbf{In-Class Activity: Instrumenting Education --- Solutions}}\\[2pt]
\end{center}

\vspace{0.5em}

A labor economist wants to estimate the \textbf{causal effect of years of education on log hourly wages}. She uses data on 3{,}010 U.S.\ men born between 1930 and 1939 from the 1980 Census. She is worried that OLS estimates are biased because \textbf{ability} is unobserved.

Proposed instrument: \textbf{quarter of birth}. Compulsory schooling laws require students to stay in school until age 16; students born earlier in the year reach 16 earlier in the school year and can legally drop out with less total schooling.

\textit{(Stylized example in the spirit of Angrist and Krueger, 1991 --- numbers constructed to prove specific pedagogical points.)}

\section*{Regression Output}

\begin{center}
{\setlength{\tabcolsep}{6pt}\renewcommand{\arraystretch}{1.05}
\begin{tabular}{lcccc}
\toprule
                   & (1)              & (2)               & (3)               & (4) \\
Dep.\ variable:    & log(wage)        & educ              & log(wage)         & log(wage) \\
Method:            & OLS              & First stage       & Reduced form      & 2SLS \\
\midrule
Years of education & 0.0700$^{***}$   &                   &                   & 0.142$^{**}$ \\
                   & (0.0035)         &                   &                   & (0.061) \\
\addlinespace
Born Q1 (Jan--Mar) &                  & $-$0.152$^{**}$   & $-$0.0216$^{**}$  & \\
                   &                  & (0.067)           & (0.0093)          & \\
\addlinespace
Black              & $-$0.236$^{***}$ & $-$1.47$^{***}$   & $-$0.342$^{***}$  & $-$0.133$^{*}$ \\
                   & (0.018)          & (0.12)            & (0.022)           & (0.075) \\
\addlinespace
Married            & 0.121$^{***}$    & 0.38$^{***}$      & 0.129$^{***}$     & 0.075$^{*}$ \\
                   & (0.016)          & (0.11)            & (0.020)           & (0.043) \\
\addlinespace
Region (South)     & $-$0.098$^{***}$ & $-$0.54$^{***}$   & $-$0.137$^{***}$  & $-$0.060 \\
                   & (0.015)          & (0.10)            & (0.018)           & (0.038) \\
\addlinespace
Constant           & 5.02$^{***}$     & 12.84$^{***}$     & 5.93$^{***}$      & 4.11$^{***}$ \\
                   & (0.052)          & (0.078)           & (0.062)           & (0.78) \\
\midrule
$N$                & 3{,}010          & 3{,}010           & 3{,}010           & 3{,}010 \\
$R^{2}$            & 0.132            & 0.087             & 0.041             & --- \\
First-stage $F$    & ---              & 5.15              & ---               & --- \\
\bottomrule
\addlinespace
\multicolumn{5}{l}{\footnotesize Standard errors in parentheses. $^{*}\,p<0.10$, $^{**}\,p<0.05$, $^{***}\,p<0.01$.} \\
\multicolumn{5}{l}{\footnotesize Column (3) is OLS of log(wage) on Born Q1 and controls. Column (4) instruments for years of education using Born Q1.} \\
\end{tabular}%
}
\end{center}

\section*{Questions + Answers}

\begin{enumerate}

\item Interpret the OLS coefficient on years of education. Be precise about units and magnitude.

\ans{One additional year of education is associated with approximately 7.0\% higher hourly wages, controlling for race, marital status, and region. (Technically: a one-year increase in education is associated with a 0.070 increase in log wages, or about 7.0\% higher wages.)}

\item Interpret the coefficient on ``Born Q1'' in the first-stage regression. What does it tell us about the relationship between quarter of birth and years of education?

\ans{Being born in Q1 (January--March) is associated with \textbf{0.152 fewer years of education} than being born in Q2--Q4, controlling for race, marital status, and region. Consistent with compulsory schooling laws: Q1 students reach the legal dropout age earlier and can exit with less total schooling.}

\item The first-stage $F$-statistic on the excluded instrument is $F = 5.15$. What does this tell us? Should we be concerned? Why or why not?

\ans{$F = 5.15$ is \textbf{below the conventional threshold of 10} for strong instruments --- a \textbf{weak instrument} concern. Consequences: (1)~2SLS estimates are biased toward OLS in finite samples; (2)~standard errors are unreliable (CIs have incorrect coverage); (3)~the 2SLS point estimate may not be trustworthy. This is a well-known limitation of the Angrist-Krueger design with a single QOB indicator. BJB (1995) formalized the critique.}

\item Interpret the 2SLS coefficient on years of education. How does it compare to the OLS estimate?

\ans{The 2SLS estimate implies that an additional year of education \textit{causes} approximately a 14.2\% increase in hourly wages --- \textbf{twice as large} as the OLS estimate of 7.0\%. Standard errors are also much larger (0.061 vs.\ 0.0035), which is typical of IV: we use less variation (only the part driven by the instrument), so estimates are noisier.}

\item With a single instrument and the same controls in all regressions, the 2SLS coefficient equals the ratio of the reduced-form coefficient on Born Q1 (column 3) to the first-stage coefficient on Born Q1 (column 2). Verify this numerically, then explain why a weak first stage is dangerous.

\ans{Numerical check: $\hat{\gamma}_{Q1}^{\text{RF}} / \hat{\pi}_{Q1}^{\text{FS}} = (-0.0216)/(-0.152) \approx 0.142$, exactly the 2SLS coefficient in column (4). So 2SLS isolates the slice of education variation that tracks with the instrument, then asks: how big is the $Y$-response per unit of that variation?

Intuition for weak-instrument danger: the denominator is small (Q1 births only shift education by 0.15 years), so any wage gap induced by Q1 gets \textit{divided by a small number}, amplifying both the point estimate \textbf{and} the noise. That's why a small first stage (low $F$) makes IV estimates large AND unreliable.}

\item Given that the OLS estimate (0.070) is \textit{smaller} than the 2SLS estimate (0.142), what does this imply about the direction of bias in OLS? Is this surprising? Why or why not?

\ans{\textbf{Surprising} if ability bias is the main problem, because ability bias should push OLS \textit{up}, not down. Three possible explanations:
\begin{enumerate}
\def\labelenumi{\arabic{enumi}.}
\item \textbf{Measurement error in education.} If education is measured with error (self-reported years), OLS is attenuated toward zero. IV corrects for classical measurement error because the instrument is correlated with \textit{true} education, not the noise. Attenuation can dominate the upward ability bias. (Leading explanation in the returns-to-schooling literature.)
\item \textbf{LATE $\neq$ ATE.} The IV estimate applies to \textit{compliers} --- marginal dropouts. Returns to schooling may be \textit{higher} for this group than the population average (heterogeneous treatment effects). So IV~$>$~OLS could reflect higher returns for compliers, not a downward OLS bias.
\item \textbf{Weak-instrument bias.} With $F = 5.15$, the 2SLS estimate is itself biased toward OLS in finite samples; the true IV (with a strong instrument) might be even larger, or the current estimate is simply unreliable.
\end{enumerate}}

\item The IV estimate is a \textbf{Local Average Treatment Effect (LATE)}. In this context, who are the \textit{compliers}?

\ans{Individuals whose total years of education were affected by compulsory schooling laws --- those who \textit{would have dropped out earlier} if they could have (i.e., if they had reached 16 before the end of the school year). These are people at the margin of the dropout decision. The IV estimate does \textbf{not} apply to always-takers (college-bound students who would have stayed regardless) or never-takers (students who dropped out before 16 regardless).}

\end{enumerate}

\end{document}
